\subsection{The integral contraction for the intermediate weights}
\label{sec:edge-contraction}

The segment calculation uses weights with only the extreme entries varying. An integral contraction compares these tests to the prescribed weights with bounded intermediate entries. At one top put \(d=c_{\rm cl}\)
and \(r=k-2\). In this subsection write \(V_\nu\) for a Schur
space in \(k\)-row matrix polynomials with column degrees
\[
 \nu_0\ge\cdots\ge\nu_{k-1}\ge0.
\]
It consists of polynomials
transforming by \(\prod_i b_{ii}^{\nu_i}\) under right
multiplication of the matrix by an upper triangular matrix \(b\).
The flag minors have column sets consisting of initial columns.
The coordinate ring of \(\mathrm{GL}_k\) under left and right
translation identifies this space in characteristic zero with the
irreducible of highest weight \(\nu\).

The products of flag minors
of this shape span that irreducible, since their span is
row-stable and contains the highest vector. The integral standard
monomial theorem gives its basis over \(O\)
\cite[Theorem II.2.16]{abw-schur}. With a diagonal term order
on the flag minors the leading monomial of a standard tableau
records the entries of its rows. These monomials are distinct
because the tableau rows are weakly increasing. Thus the
span of the basis is saturated in the entire polynomial lattice.
For a weight \(\lambda\) whose last entry may be negative put
\[
 \nu=\lambda-\lambda_{k-1}{\bf 1}.
\]

The lattice in a frame with matrix \(g\) consists of
\[
 (\det g)^{-\lambda_{k-1}}f(g^{-1}X)
\]
with \(f\) in the
standard polynomial lattice. This is the determinant twist
used in \(H(L)\).

Take end weights \(A+d,-B-d\) and intermediate weight \(\mu\)
with \(|\mu_j|\le d\), and put
\begin{equation}\label{eq:edge-intermediate-weights}\begin{aligned}
 \lambda_-&=(A,0,\ldots,0,-B),&
 \lambda_+&=(A+2d,0,\ldots,0,-B-2d),\\
 S&=(d,\mu,-d),&
 S^*&=(d,-\mu_{\rm rev},-d).
\end{aligned}\end{equation}
We use a bilinear form on the \(r\)-row Schur lattices of
middle partitions
\(\alpha=B+d+\mu\) and \(\alpha'=d-\mu_{\rm rev}\).
Their entries satisfy
\(\alpha'_i=L-\alpha_{r+1-i}\) with \(L=B+2d\).
Put \(e=\alpha_r\) and \(\gamma=\alpha-e{\bf1}\).
Division by \(\det(Y)^e\) takes the first lattice to
\(V_\gamma(O)\), since each monomial of shape \(\alpha\)
has exactly \(e\) full-size minors.

For the second lattice,
with \(J_r\) the reversal matrix, use the map
\[
             f(Y)\longmapsto (\det Y)^{L-e} f(Y^{-\mathsf T}J_r).
\]
Jacobi's complementary-minor identity sends each initial-column
minor to an initial-column complementary minor divided by
\(\det Y\), up to sign. The determinant factor clears all
denominators and the output has shape \(\gamma\).
The same transformation in reverse, by its square on invertible
matrices, gives surjectivity on the integral lattice.

\begin{proposition}[Contraction across the intermediate weights]
\label{prop:integral-edge-contraction}
At each top take the weights of \eqref{eq:edge-intermediate-weights}
with \(|\mu_i|\le d\). The flat-weight Schur lattices admit
the multiplication and contraction maps through the integral
\(r\)-row Schur lattice of shape \(\gamma\) obtained from the
two complementary partitions \(\alpha,\alpha'\). For their tensor over \(N\) tops the exponent and the corresponding \(w\)-adic order are
\[
 d_T=Nd,\qquad \operatorname{ord}_w(p^{d_T})=2d_T\quad(w^2=p,\ \operatorname{ord}_w(w)=1).
\]
 After multiplication by \(p^{d_T}\), the pulled-back target
covector satisfies the target integral test and the normalized pairing
at the distinguished edge, with a loss of at most \(2d_T\) in the
\(w\)-adic conductor at the other incoming edges. The source test
holds simultaneously at all outgoing edges.
\end{proposition}

\begin{proof}
\proofstep{A uniform pairing for the intermediate representation.}
Here is a uniform integral pairing on \(V_\gamma\).
For the transpose adjoint, the rational contravariant dual has
the same highest weight as \(V_\gamma\). Choose the corresponding
isomorphism between the two rational modules to take the highest
polynomial to the highest dual functional evaluating to one
there. Both lattices have a primitive generator on that line
and are stable under divided powers. Write \(l=\gamma_1\)
and
\[
 h=\binom r2\lfloor\log_p(l+r-1)\rfloor
\]
for \(r\ge2\).
When \(r=1\) use \(h=0\).
In either lattice the submodule generated from the highest
vector by divided powers contains \(p^h\) times the full lattice.
Indeed take a basis \(u_j\) of the divided-power lattice in
that representation.

Apply the invariant Hom projector of the
preceding type calculation to \(u_j^\vee\otimes v\), where
\(v\) lies in the full lattice. The coordinate weights of
\(V_\gamma\) lie in \([0,l]\), as each symbol appears at most
once per column of a tableau. Each coordinate of a Hom weight is a difference of two such coordinates, so
\[
 0\le a,b\le l\quad\Longrightarrow\quad -l\le a-b\le l.
\]
Multiplying by \(p^h\) gives integral intertwiners,
which take the highest vector into the integral highest line
and its divided-power span into the divided-power span.
The trace identity for that projector reconstructs \(p^h v\)
as the sum of their values on \(u_j\), multiplied by
\(\dim V_\gamma\).

Let \(M\) and \(M^\sharp\) be the two full lattices and
\(L,L^\sharp\) the respective divided-power sublattices.
The contravariant isomorphism \(j\) sends \(L\) onto \(L^\sharp\).
The inclusions just proved in both lattices give
\[
 p^{2h}M^\sharp\subseteq p^h j(M)\subseteq M^\sharp .
\]
The right inclusion uses \(p^h M\subseteq L\), and the
left uses \(p^h M^\sharp\subseteq L^\sharp\).
Thus \(p^h j\) is an integral pairing map with cokernel
killed by \(p^{2h}\). Lifting \(p^{2h}\) times a dual basis
gives the integral partners with those diagonal values. This bound depends on the width
\(l\le2d\), even though the determinant power \(e\) can be
large. In particular for \(2d\le8a\) the depth is at most
\[
 \lceil (k-1)^2\log_p(8a+k-1)\rceil.
\]

Compose with determinant division and the complementary-minor
map to obtain a pairing \(J\) on shapes \(\alpha,\alpha'\),
with the same full set of partners. Simultaneous substitution
\(Y\mapsto gY\) in both inputs multiplies \(J\) by
\(\det(g)^L\). This follows from the transpose adjoint,
since the complementary-minor input then changes by
substitution by \(g^{-\mathsf T}\) and the factor
\(\det(g)^{L-e}\). The first input contributes the remaining
\(\det(g)^e\).

\proofstep{Contraction on the polynomial coefficients.}
For a \(k\)-row polynomial \(f(X)\), put
\(D(Y)=\operatorname{diag}(1,Y,1)\).
Apply the pairing \(J\) in the variable \(Y\), coefficientwise in \(X\), to define
\[
 \mathcal C(f,t)(X)=J_Y\bigl(f(XD(Y)),t(XD(Y))\bigr).
\]

Here the two \(k\)-row inputs have the normalized shapes
of \(\lambda=\lambda_-+S\) and \(S^*\).
Cauchy--Binet expands each flag minor in these substitutions
into middle flag minors of \(Y\) with coefficients in \(O[X]\)
for integral input. The \(Y\)-shapes are \(\alpha,\alpha'\)
because the first column of \(X\) is retained and the last is
unused. Thus contraction is integral. It commutes with left matrix substitution in \(X\).

For an invertible middle block \(g\) one has
\[
 D(g)D(Y)=D(gY),
\]
so the covariance of \(J\) gives
\[
 \mathcal C(f,t)(XD(g))=\det(g)^L\,\mathcal C(f,t)(X).
\]
This proves the middle diagonal degrees and invariance under
upper unipotent middle-column operations. A first-column
addition to a middle column is represented by a matrix \(u\)
whose only off-diagonal entry is in the first row. The matrix
\(D(Y)^{-1}uD(Y)\) is upper unipotent over \(K(Y)\), so
\[
 XuD(Y)=XD(Y)(D(Y)^{-1}uD(Y))
\]
has the same value in each input. Invariance there gives a
polynomial identity over \(K\). The last column is unused,
and the first-column degrees add. Together these operations
give precisely the right triangular character of the
normalized shape of \(\lambda_+\).

For any middle polynomial of shape \(d+\mu\), lift each
flag minor of size \(j\) to a flag minor including the first
row and first column, of size \(j+1\). If the product has \(a\) minors, the lifted product has first-row degree \(a\). Multiply by
\[
 X_{00}^{2d-a},\qquad 0\le a\le2d,
\]
to obtain first-row degree \(2d\).
Evaluation on \(D(Y)\) gives the original polynomial, and
the lift has shape \(S+d{\bf1}\) and last-row degree zero.
The same construction lifts the weight \(d-\mu_{\rm rev}\)
to the shape of \(S^*\). Since multiplication by
\((\det Y)^B\) identifies the shape \(d+\mu\) with \(\alpha\),
choose the lifted partners \(s_b,t_a\) so that
\[
 J((\det Y)^B s_b(D(Y)),t_a(D(Y)))
\]
is the chosen
scalar diagonal.

If \(F\) has the lower flat shape, its
substitution \(F(XD(Y))\) is \((\det Y)^B F(X)\),
by the first-column and cofactor formula.
It follows that \(\mathcal C(Fs_b,t_a)\) equals \(F\)
times a polynomial with the flat normalized shape
\((4d,2d,\ldots,2d,0)\), first-row degree \(4d\)
and last-row degree zero. This space is spanned by
\[
          Z=(X_{00}\det X_{[0,k-2],[0,k-2]})^{2d}.
\]
In its first-column and cofactor model, every variable factor
must have maximal first-row degree, and each cofactor must
omit the last row. Since \(Z(I_k)=1\), evaluation at the
identity matrix shows that
\[
 \mathcal C(Fs_b,t_a)= \delta_{ab}p^{2h}ZF.
\]

This constructs the one-component maps
\[
 i(F)=ZF,\ U_b(F)=s_b F,\ L_a(G)=\mathcal C(G,t_a).
\]

The number of pairs is \(\dim V_\gamma\), the dimension for
the middle weight. Tensor over all tops, reversing the two
end degrees at the odd tops, to obtain the maps on a cycle.

\proofstep{Verification on the edge lattices.}
We check the edge lattices for those maps. In a frame of an
edge the single changing row direction in the lattice has
scale \(1,w^{-1},w^{-2}\) at the source, intermediate
lattice and target. For first-row degree \(2d\) and last weight \(-d\) the determinant-twisted scaling factors are
\[
 w^d\quad\text{at the first change},\qquad w^{2d}\quad\text{at the second}.
\]
Thus the transported
\(s_b,t_a\) belong to the normalized intermediate lattice,
and \(w^{2d}t_a\) is integral at the target.
For \(Z\) the degree and the last weight are doubled.
It is primitive in the coefficient lattice of the intermediate
frame after normalization. To check saturation, reduce a product \(ZP\) in the coefficient
lattice modulo the uniformizer \(w\). Primitivity of \(Z\) and the
domain property give
\[
 \overline Z\,\overline P=0,\quad\overline Z\ne0
       \ \Longrightarrow\ \overline P=0 .
\]
The original coefficient lattices are saturated, so \(P/w\) then
belongs to the smaller lattice. This proves saturation of the image.

For any outgoing edge at the same source, a polynomial
integral at that source belongs to its normalized edge
lattice. Indeed in the respective adapted row coordinates
the transformed monomials have row weights between the
two end weights. In particular the normalization subtracts
the greatest one. This proves the outgoing tests for
\(i\) and \(U_b\), using integrality of multiplication
in every frame. For any incoming edge at the same target,
the target lattice for \(S^*\) likewise lies in its
normalized intermediate lattice. In the edge frame this uses
\[
 2j-(j-d)=j+d\ge0
\]
for every row weight \(j\in[-d,d]\).
For the tensor of \(t_a\)'s put \(d_T=Nd\).

At the target \(\mathcal C(\,\cdot\,,t_a)\) is integral after
multiplication by \(w^{2d_T}\), and it is integral at the
distinguished normalized edge with no multiplier.
At any other incoming normalized edge it loses at most
\(w^{2d_T}\), by the target inclusion just proved.
The normalized first weights of \(\lambda_+\) exceed those
of \(\lambda\) by \(d_T\). The pullback of a covector
multiplied by \(p^{d_T}\) therefore preserves the target
condition and the raw edge pairing, and decreases the
vanishing conductor at other incoming edges by at most
\(2d_T\). These are the properties needed in the next
subsection. They use the same integral section at all
neighbors of the fixed source or target.
\end{proof}
